\documentclass{SCNU_Experimental_Report}

% 填写报告头信息
\studentname{Tommy}
\studentid{20240110}
\major{电子信息工程}
\class{24级01班}
\coursename{数字图像处理}
\projectname{图像处理的基本操作和图像直方图均衡化}
\verify
\experimentdate{2024年10月1日}
\teacher{Tommy}

\begin{document}
\maketitle

\section{实验目的}
\begin{enumerate}
    \item 掌握图像的离散傅里叶变换的幅度谱和相位谱；
    \item 通过实验了解二维频谱的分布特点；
    \item 熟练掌握 FFT 的变换方法及应用；
    \item 熟练掌握傅里叶变换的基本性质。
\end{enumerate}

\section{实验平台}
装有 MATLAB 软件的计算机一台

\section{实验内容}
\begin{enumerate}
    \item 显示图像的离散傅里叶变换的幅度谱和相位谱；
    \item 二维离散傅里叶逆变换；
    \item 傅里叶变换谱实部的二维离散傅里叶逆变换；
    \item 傅里叶变换相位信息的二维离散傅里叶逆变换；
    \item 图像二维离散傅里叶变换的性质。
\end{enumerate}

\section{实验原理}
当数字图像 \(f(x,y)\) 为 \(N \times N\) 个像素时，其对应的二维离散傅里叶变换为：
\begin{equation}
    F(u,v) = \frac{1}{N}\sum_{x=0}^{N-1}\sum_{y=0}^{N-1} f(x,y) \exp\left[-\frac{\mathrm{j}2\pi(ux+vy)}{N}\right], \quad u,v = 0,1,\ldots,N-1
\end{equation}
从图像的二维离散傅里叶变换幅度谱中可见：一般图像的频谱大部分集中在直流和低频区域，而高频成分相对分量较低，高频信息对应着图像的细节或边界或突变地方或噪声；图像低频信息对应着变化缓慢区域。在数字图像中，相位信息非常重要，因为相位信息中携带着图像的边界位置等结构信息。
二维离散傅里叶逆变换为：
\begin{equation}
    f(x,y) = \frac{1}{N}\sum_{u=0}^{N-1}\sum_{v=0}^{N-1} F(u,v) \exp\left[\frac{\mathrm{j}2\pi(ux+vy)}{N}\right]
\end{equation}
\subsection{二维离散傅里叶变换的性质}
若傅里叶变换对表示如下：\(f(x,y) \Leftrightarrow F(u,v)\)

\section{实验步骤}
\subsection{显示图像的离散傅里叶变换的幅度谱和相位谱}
\begin{figure}[H]
    \centering
    \includegraphics[width=0.4\textwidth]{fig1.png}
    \caption{原图}
\end{figure}

\section{思考题}
\subsection{二维离散傅里叶的可分离性有什么意义？}
答：二维离散傅里叶变换的可分离性意味着可以将二维 DFT 分解为两次一维 DFT 的连续计算（先对每一行进行一维 DFT，再对每一列进行一维 DFT）。其意义在于计算效率大幅提升：直接计算二维 DFT 的复杂度为 \(O(N^4)\)，利用可分离性可降至 \(O(2N^3)\)，而 FFT 可进一步降至 \(O(N^2 \log N)\)。

\section{基础三线表}
\begin{table}[H]
    \centering
    \caption{MATLAB 数据类型对照表}
    \label{tab:datatype}
    \begin{tabular}{clc}
        \toprule
        类型 & 描述 & 范围 \\
        \midrule
        uint8 & 无符号8位整数 & 0 -- 255 \\
        uint16 & 无符号16位整数 & 0 -- 65535 \\
        int8 & 有符号8位整数 & -128 -- 127 \\
        double & 双精度浮点数 & $-10^{308}$ -- $10^{308}$ \\
        logical & 逻辑值 & 0 或 1 \\
        \bottomrule
    \end{tabular}
\end{table}

\end{document}